\chapter{Entropy, and the Cost of Staying Level}
\label{ch:homeo-ent}
% ===================================================================

The previous chapter gave the learner a goal it can state about itself,
and a reading of its own distance from that goal.
It did not say what the goal costs.

This chapter says.
Homeostasis is not free, and it is not free for a reason that has been
sitting in Part~\ref{part:physics} since the conservation chapter: a
cost-accounted theory is frictionful by construction, and a learner
holding a level against friction is paying a bill in exactly the currency
that chapter priced.
What follows is short, because most of the work was already done
elsewhere.
Its job is to connect it.

\section{Dissipation is where information leaves}

Chapter~\ref{sec:cost} established that a cost-accounted theory is
frictionful by design and therefore lives in the dissipative regime:
every lossy enzyme that fires strictly decreases the free energy, and the
virtual token is a Lyapunov function rather than a conserved charge.

Read that result backwards.
A frictionless theory is one whose enzymes are exactly value-preserving,
and by Theorem~\ref{thm:vtok-ftap} such a theory has an exact potential,
which is to say it has enough structure to say where every unit of value
came from.
A frictionful theory does not.
The dissipation $\theta$ accumulated around a cycle is precisely the
amount of provenance the theory cannot reconstruct: value went somewhere,
and the ledger does not record where.

That is an information-theoretic statement wearing economic clothing, and
Chapter~\ref{sec:cost} already turns it into one, by way of the Landauer
bound: the erasure a computation performs is bounded below by the
holonomy of the energy form.
Erasure is forgetting which of several possible predecessors a state
had.
Dissipation is the price of forgetting.
The two thresholds coincide, which is the most satisfying result in that
chapter and is also, for present purposes, a signpost.

\section{What a setpoint costs}

Now put the bathtub of Chapter~\ref{ch:suffer} in that setting.

A learner holding $\vect{A}$ near $\setpt$ is not sitting still.
It is running an assay on its own account, comparing the result to a held
formula, and adjusting an outflow it does not fully control against an
inflow it does not control at all.
Every one of those operations is a metered rewrite, and by
Definition~\ref{def:affordable} every metered rewrite is drawn from the
same account it is measuring.

\begin{proposition}[Regulation is charged to what it regulates]
\label{ent:prop:self-charge}
Let $\Lrn$ be an internalizing learner in the sense of
Definition~\ref{suf:def:internal}, and let $\kappa_{\mathrm{reg}}$ be the
per-cycle cost of the assay, the comparison, and the corrective rewrite.
Then $\kappa_{\mathrm{reg}}$ is charged to $\vect{A}$ itself.
Consequently the inflow required to hold $\vect{A} = \setpt$ is strictly
greater, by $\kappa_{\mathrm{reg}}$ per cycle, than the inflow required to
sustain a non-internalizing learner with the same rewrite load.
\end{proposition}

\begin{proof}
The account is vectorial and single: there is no second account from
which regulation could be funded, since by Remark~\ref{rmk:ambient} an
account merely adjacent to a process is ambient authority, which the
framework has already refused.
So the assay of \S\ref{suf:S3} draws on $\vect{A}$, and the drift it
reports is a drift of the quantity that just paid for the report.
The stated inequality is the difference of the two flow balances.
\end{proof}

\begin{remark}[The thermostat pays for the thermometer]
\label{ent:rmk:thermostat}
Proposition~\ref{ent:prop:self-charge} is the reason the previous chapter
could not simply assert that internalizing is better.
A learner that watches its own reservoir is, by that act, draining it.
The watching is worth doing only when what it buys --- anticipation,
which is to say correction begun before the level has actually fallen ---
exceeds what it costs, and whether it does is a fact about the
environment rather than about the learner.

This is also the sharpest available statement of what suffering is
\emph{for}, and it is unsentimental.
Suffering is the readout of a meter the organism is paying to run.
An organism that felt nothing would be cheaper, and in a placid world it
would win.
\end{remark}

\section{The two erasures, and what is left over}

\begin{remark}[Two erasures, and the gap between them]
\label{rmk:two-erasures}
There are in fact two distinct erasures in play, and they are not the
same map~\cite{TwoErasures}.
One forgets the \emph{history}: the history monad's counit, which
discards the log and keeps the state.
The other forgets the \emph{cost}: the cost monad's counit, which
discards the meter and keeps the computation.
Neither factors through the other, and the failure to factor is graded
--- there is a lattice of intermediate laxities between them, each
corresponding to a theory that remembers some of one and some of the
other.

The gap between the two erasures is not a technical annoyance.
It is a supply of positions.
A theory sitting strictly between them knows something about its own past
that it cannot express as a cost and something about its own costs that
it cannot express as a history, and a computation with information of
that shape is exactly a computation with something to decide.
\end{remark}

An internalizing learner sits in that gap by construction, which is worth
saying plainly because it was not obvious that anything did.

Its held setpoint is cost information: $\setpt$ is a level in the
account, and $\phi_{\mathrm{home}}$ is a formula about the meter.
Its drift is history information: a drift is a difference taken across
time, and reading one requires having kept what the level was.
Neither reduces to the other.
A learner that kept only the cost would know its level and not know
whether it was falling; a learner that kept only the history would know
the shape of its trace and not what it could afford.
Suffering, as Chapter~\ref{ch:suffer} defines it, is precisely the datum
that requires both --- which is why it could not be a component of the
account, and had to be a formula about one.

\begin{remark}[Where this goes]
\label{ent:rmk:handoff}
That leaves an obligation this part does not discharge.
If an internalizing learner is a computation with something to decide,
then something must do the deciding, and nothing constructed so far says
what.
The reduction relation offers continuations and is silent about which
occurs.

The Prestige takes that silence seriously, and argues that it is where
the only genuinely non-Turing thing in this book is located.
Readers who want that argument now will find it in
Chapter~\ref{ch:ent}.
Readers content to let the Turn finish its own business should carry
forward only this: the learner of these chapters has been given, for the
first time, information of a shape that makes a choice meaningful, and it
was homeostasis that gave it.
\end{remark}

% ===================================================================
